\section{The Service That Owns the Speakers}
\label{sec:ch09-the-service-that-owns-the-speakers}

\index{Speakers service!the project file|(}%
A new folder beside the first one, and a project file that is
chapter~\ref{ch:first-subgraph}'s with a different name on it:

\begin{minted}[fontsize=\footnotesize]{xml}
<Project Sdk="Microsoft.NET.Sdk.Web">

  <PropertyGroup>
    <TargetFramework>net10.0</TargetFramework>
    <Nullable>enable</Nullable>
    <ImplicitUsings>enable</ImplicitUsings>
    <GenerateDocumentationFile>true</GenerateDocumentationFile>
    <NoWarn>$(NoWarn);CS1591</NoWarn>
  </PropertyGroup>

  <ItemGroup>
    <PackageReference Include="HotChocolate.ApolloFederation" Version="16.6.1" />
    <PackageReference Include="HotChocolate.AspNetCore" Version="16.6.1" />
    <PackageReference Include="HotChocolate.AspNetCore.CommandLine" Version="16.6.1" />
    <PackageReference Include="HotChocolate.Data.EntityFramework" Version="16.6.1" />
    <PackageReference Include="HotChocolate.Types.Analyzers" Version="16.6.1" />
    <PackageReference Include="Microsoft.EntityFrameworkCore.Sqlite" Version="10.0.11" />
  </ItemGroup>

</Project>
\end{minted}

\index{ports!5002 is pinned in a launch profile}%
The launch profile pins the port, for the reason
chapter~\ref{ch:hot-chocolate-zero} pinned 5001: \code{dotnet new web} picks
one at random, and a book printing raw HTTP would otherwise print fiction.

\begin{minted}{json}
{
  "$schema": "https://json.schemastore.org/launchsettings.json",
  "profiles": {
    "http": {
      "commandName": "Project",
      "dotnetRunMessages": true,
      "launchBrowser": false,
      "applicationUrl": "http://localhost:5002",
      "environmentVariables": {
        "ASPNETCORE_ENVIRONMENT": "Development"
      }
    }
  }
}
\end{minted}

\subsection{The Type, and the Table Under It}

\index{Speakers service!declares its own Speaker}%
\code{Speaker} is the record chapter~\ref{ch:hot-chocolate-zero} wrote, in a
new namespace and with one sentence added to say which service holds the rows
now. The three properties are untouched:

\begin{minted}{csharp}
namespace Speakers;

/// <summary>
/// A person giving one or more sessions. This service owns the rows;
/// the Sessions subgraph knows only that a speaker with this id exists.
/// </summary>
/// <param name="Id">Stable identifier for this speaker.</param>
/// <param name="Name">The name to print in the program.</param>
/// <param name="Bio">A sentence or two, if the speaker sent any.</param>
public sealed record Speaker(int Id, string Name, string? Bio);
\end{minted}

\index{Speakers service!its own database file}%
Its own database, in its own file, seeded from its own data. Three speakers,
the same three the conference has had since chapter~\ref{ch:hot-chocolate-zero}:

\begin{minted}{csharp}
namespace Speakers;

/// <summary>
/// The speaker list, seeded into this service's own database. These are
/// the rows chapter 3 seeded into the Sessions service; they live here
/// now, and only here.
/// </summary>
public static class SpeakerData
{
    public static readonly IReadOnlyList<Speaker> Speakers =
    [
        new Speaker(1, "Ada Fischer", "Works on query planners."),
        new Speaker(2, "Bruno Kaminski", null),
        new Speaker(3, "Chidi Okafor", "Runs platform engineering.")
    ];
}
\end{minted}

\begin{minted}{csharp}
using Microsoft.EntityFrameworkCore;

namespace Speakers;

/// <summary>
/// This service's own database, in its own file. Nothing here knows
/// that sessions exist: a speaker does not carry a list of them, and
/// no query in this service joins across that seam because there is
/// nothing on the other side of it to join to.
/// </summary>
public sealed class SpeakerContext(
    DbContextOptions<SpeakerContext> options) : DbContext(options)
{
    public DbSet<Speaker> Speakers => Set<Speaker>();

    protected override void OnModelCreating(ModelBuilder model)
        => model.Entity<Speaker>().HasData(SpeakerData.Speakers);
}
\end{minted}

\index{seam!a speaker does not know its sessions}%
Read what is missing from that context, because it is the seam seen from the
other side. A speaker has no list of sessions. It could not have one: the
sessions are in a different database in a different process, and no
\code{DbSet} reaches them. That absence is not a simplification made for the
book: any query that needs a speaker and a session in the same answer now needs
two services to produce it, and no amount of care inside this one changes that.
Chapter~\ref{ch:cross-seam-paging} is a whole chapter about what it costs.

\subsection{Two Files Typed Out Again}

\index{DataLoader!each service declares its own}%
The loader is chapter~\ref{ch:data-without-n-plus-1}'s, against this service's
context:

\begin{minted}{csharp}
using Microsoft.EntityFrameworkCore;

namespace Speakers;

public static class SpeakerDataLoader
{
    /// <summary>
    /// Every speaker one batch asked for, keyed by identifier. The
    /// router sends this service a list of representations rather than
    /// one at a time, so the loader behind the reference resolver is
    /// what decides whether that list costs one statement or several.
    /// </summary>
    [DataLoader]
    public static async Task<Dictionary<int, Speaker>>
        GetSpeakerByIdAsync(
            IReadOnlyList<int> ids,
            SpeakerContext db,
            CancellationToken ct)
        => await db.Speakers
            .Where(s => ids.Contains(s.Id))
            .ToDictionaryAsync(s => s.Id, ct);
}
\end{minted}

\index{ShareNodeFields@\code{ShareNodeFields}!declared again in this service}%
And chapter~\ref{ch:composition}'s type interceptor. Every line from the class
declaration down is the same as the Sessions copy; only the namespace and the
comment explaining why there are two of it differ:

\begin{minted}{csharp}
using HotChocolate.Configuration;
using HotChocolate.Language;
using HotChocolate.Types.Descriptors.Configurations;

namespace Speakers;

/// <summary>
/// Applies @shareable to the two fields global object identification
/// generates, and to nothing else. This is chapter 7's file, declared
/// again in this service rather than shared with the Sessions one:
/// every subgraph in this graph exports node and nodes, so every
/// subgraph needs this, and a project reference between two services
/// that would be separate repositories is the thing this book is
/// arguing against.
/// </summary>
internal sealed class ShareNodeFields : TypeInterceptor
{
    public override void OnBeforeCompleteType(
        ITypeCompletionContext context,
        TypeSystemConfiguration configuration)
    {
        if (configuration is not ObjectTypeConfiguration type
            || type.Name != "Query")
        {
            return;
        }

        foreach (var field in type.Fields)
        {
            if (field.Name is "node" or "nodes")
            {
                field.Directives.Add(
                    new DirectiveConfiguration(
                        new DirectiveNode("shareable")));
            }
        }
    }
}
\end{minted}

\index{duplication!is the price of no shared library}%
That is the second copy of a class in a book that spends a chapter on why
duplication kills schemas, and I want to be plain about the trade rather than
let it pass. It exists twice because the alternative is a package or a project
reference between two services that are meant to be independently deployable. A shared library between subgraphs is a coupling
that does not show up in the schema, does not show up in composition, and
shows up the first time two teams need different versions of it on different
release trains. I would rather type it twice. \code{StatementLog} from
chapter~\ref{ch:data-without-n-plus-1} is here as well, for the same reason and
by the same rule, and it is identical apart from its namespace.

\subsection{The Entity, and the Root Field}

\index{Speakers service!carries chapter 6's SpeakerType}%
\code{SpeakerType} is chapter~\ref{ch:first-subgraph}'s file. It has moved
service and changed nothing that matters:

\begin{minted}{csharp}
using HotChocolate.ApolloFederation.Resolvers;
using HotChocolate.ApolloFederation.Types;
using HotChocolate.Types.Relay;

namespace Speakers;

[ObjectType<Speaker>]
[Key("id")]
[ReferenceResolver(
    EntityResolver = nameof(ResolveSpeakerReferenceAsync))]
public static partial class SpeakerType
{
    /// <summary>
    /// Loads one speaker by the identifier inside a global node id.
    /// </summary>
    [NodeResolver]
    public static async Task<Speaker?> GetSpeakerAsync(
        int id,
        ISpeakerByIdDataLoader speakerById,
        CancellationToken ct)
        => await speakerById.LoadAsync(id, ct);

    /// <summary>
    /// Answers one representation from the entity route. This is the
    /// method the router reaches every time a client asks a session
    /// for its speaker's name, so it is the busiest resolver in the
    /// graph. The key arrives as the undecoded node id string, and the
    /// type name inside it is checked before the integer is spent.
    /// </summary>
    [GraphQLIgnore]
    public static async Task<Speaker?> ResolveSpeakerReferenceAsync(
        string id,
        INodeIdSerializer serializer,
        ISpeakerByIdDataLoader speakerById,
        CancellationToken ct)
    {
        var nodeId = serializer.Parse(id, typeof(int));

        if (nodeId.TypeName != nameof(Speaker))
        {
            return null;
        }

        return await speakerById.LoadAsync(
            (int)nodeId.InternalId!, ct);
    }
}
\end{minted}

\index{reference resolver!becomes the busiest method in the graph}%
One comment in it says something chapter~\ref{ch:first-subgraph}'s could not.
There, this reference resolver was a route nobody used: the entity route
existed, it answered, and no router asked it anything. From this chapter on it
is the method every cross-seam request in the book goes through, which makes
the DataLoader behind it the difference between two statements and a number
that grows with the page.

\index{Query@\code{Query}!this service gets a root field of its own}%
A subgraph with no root field of its own is legal and strange, so this one has
one:

\begin{minted}{csharp}
using GreenDonut.Data;

namespace Speakers;

[QueryType]
public static class Query
{
    /// <summary>
    /// Everyone speaking at the conference, in program order.
    /// </summary>
    [UsePaging(DefaultPageSize = 2, MaxPageSize = 10)]
    public static IQueryable<Speaker> GetSpeakers(
        SpeakerContext db,
        QueryContext<Speaker> query)
        => db.Speakers
            .OrderBy(s => s.Name)
            .With(query);
}
\end{minted}

\index{PageInfo@\code{PageInfo}!two subgraphs now generate it}%
Paging that field makes this the second subgraph to generate a \code{PageInfo}
type, and the two would collide on it if the package had not already dealt with
that. Chapter~\ref{ch:composition} found \code{@shareable} written onto all four
of that type's fields without being asked, and called it an oddity when one
service was doing it. Here is what it was for: two services generate the same
type, both declare every field of it shareable, and the composition nobody
worried about passes.

\subsection{Program.cs, and What Is Not In It}

\begin{minted}{csharp}
using ChilliCream.Nitro.App;
using HotChocolate.AspNetCore;
using Microsoft.EntityFrameworkCore;
using Speakers;

var builder = WebApplication.CreateBuilder(args);

// EF Core reports every command through the logging pipeline as well,
// at Information and with a duration attached. StatementLog below is
// this service's instrument, so silence the other one rather than have
// every statement printed twice.
builder.Logging.AddFilter(
    "Microsoft.EntityFrameworkCore.Database.Command",
    LogLevel.Warning);

builder.Services.AddDbContextFactory<SpeakerContext>(options =>
{
    options.UseSqlite("Data Source=speakers.db");

    if (builder.Environment.IsDevelopment())
    {
        options.AddInterceptors(new StatementLog());
    }
});

builder
    .AddGraphQL()
    .AddSpeakersTypes()
    .RegisterDbContextFactory<SpeakerContext>()
    .AddQueryContext()
    .AddPagingArguments()
    .AddGlobalObjectIdentification()
    .AddApolloFederation()
    .ModifyCostOptions(options => options.ApplyCostDefaults = false)
    .TryAddTypeInterceptor<ShareNodeFields>();

var app = builder.Build();

using (var scope = app.Services.CreateScope())
{
    var factory = scope.ServiceProvider
        .GetRequiredService<IDbContextFactory<SpeakerContext>>();
    using var db = factory.CreateDbContext();
    db.Database.EnsureCreated();
}

app.MapGraphQL()
    .WithOptions(options =>
    {
        // Serve the Nitro build that shipped in the package. The
        // default, ServeMode.Latest, downloads the current one
        // from ChilliCream's CDN instead.
        options.Tool.ServeMode = ServeMode.Embedded;
    });

app.RunWithGraphQLCommands(args);
\end{minted}

\index{AddMutationConventions@\code{AddMutationConventions}!absent here}%
One line from the Sessions builder is missing and its absence is not an
oversight. The call to
\code{AddMutationConventions()}
generates the union and the error interface around a mutation, and this
service has no mutation. Nothing else differs: the cost option
chapter~\ref{ch:composition} needed to get past the composer is here, because
this service would emit the same unasked-for \code{@cost} weights without it.

\index{Speakers service!answers on its own port}%
Start it and ask it something, on 5002, with no router involved:

\begin{minted}{text}
POST /graphql HTTP/1.1
Host: localhost:5002
Content-Type: application/json

{"query":"{ speakers { nodes { id name } } }"}
\end{minted}

\begin{minted}[fontsize=\footnotesize,breakanywhere]{json}
{"data":{"speakers":{"nodes":[{"id":"U3BlYWtlcjox","name":"Ada Fischer"},{"id":"U3BlYWtlcjoy","name":"Bruno Kaminski"}]}}}
\end{minted}

Two speakers rather than three. The page size comes from the attribute on
\code{GetSpeakers} above, set to the same default
chapter~\ref{ch:data-without-n-plus-1} chose for sessions.

\index{node id!is the same string in both services}%
The line to read is the first id. \code{U3BlYWtlcjox} is the string chapter~\ref{ch:schema-design} produced
in the Sessions service for the same speaker, out of a different process
against a different database, and nothing coordinated the two. Both services
encode a type name and a primary key the same way because both use the same
default, and that is the whole reason the seam in the next section works at
all.
\index{Speakers service!the project file|)}%
